% ============================================================
% EUROPASS CURRICULUM VITAE
% Shamekh Al-Suwi
%
% Format: Europass CV
% Language: English (UK)
% Compiler: pdfLaTeX
%
% Required LaTeX packages:
%   europasscv
%   babel
%   hyperref (loaded by europasscv)
%
% Updated: October 2026
% ============================================================

% !TEX encoding = UTF-8
% !TEX program = pdflatex
% !TEX spellcheck = en_GB

\documentclass[english,a4paper]{europasscv}

\usepackage[english]{babel}

% ============================================================
% PERSONAL INFORMATION
% ============================================================

\ecvname{Shamekh Al-Suwi}

\ecvaddress{4020 Linz, Upper Austria, Austria}

\ecvmobile{+43 677 63825795}

\ecvemail{shamekhing@gmail.com}

\ecvhomepage{www.linkedin.com/in/shamekhing}

\ecvdateofbirth{23/04/1996}

\ecvnationality{Jordanian}

% Optional photograph:
% \ecvpicture[width=3.8cm]{picture.jpg}

% ============================================================
% DOCUMENT
% ============================================================

\begin{document}

\begin{europasscv}

% ============================================================
% PROFILE
% ============================================================

\ecvpersonalinfo

\ecvbigitem{About me}{
Multidisciplinary AI and software engineer with a background
in mechatronics, control systems, computer vision, and
applied research.

Professional experience spanning artificial intelligence,
industrial software development, scientific computing,
automation, and research engineering.

Currently pursuing an MSc in Artificial Intelligence at
Johannes Kepler University Linz, with additional independent
work in computer architecture, FPGA development,
SystemVerilog, and hardware acceleration.

Research interests include efficient machine learning
accelerators, hardware-software co-design, embedded
intelligence, and the application of large language models
to hardware design and verification.
}

\ecvbigitem{Job applied for}{
AI Research Engineer / Embedded Systems Engineer /
FPGA Engineer / Software Engineer
}

% ============================================================
% WORK EXPERIENCE
% ============================================================

\ecvsection{Work experience}

% ------------------------------------------------------------
% FREELANCE AND INDEPENDENT DEVELOPMENT
% ------------------------------------------------------------

\ecvtitle{March 2026 -- Present}{
Freelance Software Developer
}

\ecvitem{}{
Independent software development and technical projects
}

\ecvitem{}{
\begin{ecvitemize}

    \item Development of C++ applications using
    C++17, CMake, and Raylib.

    \item Development of Python-based tools,
    automation utilities, and computer vision applications.

    \item Application of large language models to
    software development workflows.

    \item Independent exploration of embedded systems,
    digital logic, and FPGA-based computing architectures.

    \item Development and experimentation with
    simulation and verification workflows on Linux.

\end{ecvitemize}
}

\ecvitem{}{
\ecvhighlight{Business or sector}
\quad Software Development / Independent Engineering
}

% ------------------------------------------------------------
% AUSTRIAN INSTITUTE OF TECHNOLOGY
% ------------------------------------------------------------

\ecvtitle{October 2024 -- March 2026}{
Junior Research Engineer
}

\ecvitem{}{
\href{https://www.ait.ac.at/en/}{
Austrian Institute of Technology (AIT)
}
-- Technology Experience Department
}

\ecvitem{}{
\begin{ecvitemize}

    \item AI system integration and machine learning
    system development for production environments.

    \item Development of data-driven monitoring systems
    and real-time AI interfaces.

    \item Prototype-to-deployment workflows and
    research-to-application transfer.

    \item Applied research in system reliability,
    evaluation, and user experience.

    \item Integration of software components
    in experimental research environments.

\end{ecvitemize}
}

\ecvitem{}{
\ecvhighlight{Business or sector}
\quad IT / Applied Research
}

% ------------------------------------------------------------
% JKU COMPUTER GRAPHICS
% ------------------------------------------------------------

\ecvtitle{October 2023 -- February 2024}{
Teaching Assistant
}

\ecvitem{}{
\href{https://www.jku.at/institut-fuer-computergrafik/}{
Johannes Kepler University Linz
}
-- Institute of Computer Graphics
}

\ecvitem{}{
\begin{ecvitemize}

    \item Supported practical exercises in computer vision,
    machine learning, and deep learning.

    \item Assisted students with implementations,
    data analysis, and system modelling.

    \item Supported course materials and
    practical learning activities.

\end{ecvitemize}
}

\ecvitem{}{
\ecvhighlight{Business or sector}
\quad Higher Education / Academia
}

% ------------------------------------------------------------
% SEQUALITY
% ------------------------------------------------------------

\ecvtitle{October 2023 -- November 2023}{
Software Developer Intern
}

\ecvitem{}{
\href{
https://www.sequality.at/en/your-partner-for-technical-software/
}{
Sequality
}
}

\ecvitem{}{
\begin{ecvitemize}

    \item Developed industrial software user interfaces
    using C++ and Qt.

    \item Contributed to technical software applications
    for industrial environments.

    \item Worked with performance-sensitive software
    and reliable system interfaces.

\end{ecvitemize}
}

\ecvitem{}{
\ecvhighlight{Business or sector}
\quad Industrial Software Engineering
}

% ------------------------------------------------------------
% JKU SEMICONDUCTOR PHYSICS
% ------------------------------------------------------------

\ecvtitle{August 2022 -- February 2023}{
Research Assistant
}

\ecvitem{}{
\href{
https://www.jku.at/en/institute-of-semiconductor-and-solid-state-physics/
}{
Johannes Kepler University Linz
}
-- Institute of Semiconductor and Solid State Physics
}

\ecvitem{}{
\begin{ecvitemize}

    \item Scientific software maintenance for a
    quantum key distribution repository.

    \item Worked with software relating to
    entangled-photon quantum key distribution.

    \item Maintained components supporting secure
    distributed systems and cryptographic research.

    \item Contributed to research-grade software
    repositories and scientific computing workflows.

\end{ecvitemize}
}

\ecvitem{}{
\ecvhighlight{Business or sector}
\quad Scientific Research / Higher Education
}

% ------------------------------------------------------------
% SIMMA SOFTWARE
% ------------------------------------------------------------

\ecvtitle{November 2021 -- April 2022}{
Software Engineer
}

\ecvitem{}{
\href{https://www.simma.io/}{
Simma Software
}
}

\ecvitem{}{
\begin{ecvitemize}

    \item Developed automation software
    using Selenium.

    \item Implemented software automation pipelines
    and system integration workflows.

    \item Developed tools supporting
    workflow reliability and automated testing.

\end{ecvitemize}
}

\ecvitem{}{
\ecvhighlight{Business or sector}
\quad Information Technology
}

% ------------------------------------------------------------
% MIXED DIMENSIONS
% ------------------------------------------------------------

\ecvtitle{February 2020 -- November 2021}{
R\&D Software Engineer
}

\ecvitem{}{
\href{https://integralrealitylabs.com/}{
Mixed Dimensions
}
}

\ecvitem{}{
\begin{ecvitemize}

    \item Research and development of geometrical
    algorithms and computational models.

    \item Developed software for repairing damaged
    3D models for commercial 3D printing.

    \item Implemented geometry-processing algorithms
    and data-processing pipelines.

    \item Worked on computational efficiency
    and software performance optimisation.

\end{ecvitemize}
}

\ecvitem{}{
\ecvhighlight{Business or sector}
\quad Software Research and Development
}

% ------------------------------------------------------------
% BAHBISH
% ------------------------------------------------------------

\ecvtitle{January 2019 -- April 2019}{
Mechatronics Intern
}

\ecvitem{}{Bahbish}

\ecvitem{}{
\begin{ecvitemize}

    \item Mechatronics engineering internship.

    \item Microscope automation.

    \item Computer-aided design (CAD).

\end{ecvitemize}
}

\ecvitem{}{
\ecvhighlight{Business or sector}
\quad Engineering
}

% ============================================================
% EDUCATION
% ============================================================

\ecvsection{Education and training}

% ------------------------------------------------------------
% MASTER'S DEGREE
% ------------------------------------------------------------

\ecvtitlelevel{
March 2024 -- Present
}{
MSc Artificial Intelligence
}{
ISCED~7
}

\ecvitem{}{
\href{https://www.jku.at/en/}{
Johannes Kepler University Linz
}, Austria
}

\ecvitem{Study Areas}{
\begin{ecvitemize}

    \item Machine learning and deep learning.

    \item Neural networks and recurrent architectures.

    \item Deep reinforcement learning.

    \item Computer vision.

    \item Probabilistic models and statistical learning.

    \item Parallel computing.

    \item Signal processing and adaptive systems.

    \item Quantum information.

\end{ecvitemize}
}

\ecvitem{Relevant Coursework}{
\begin{ecvitemize}

    \item \textbf{Machine Learning:}
    Supervised Techniques, Unsupervised Techniques,
    Deep Learning I and II.

    \item \textbf{Neural Networks:}
    LSTM and Recurrent Neural Networks,
    Deep Reinforcement Learning.

    \item \textbf{Mathematics and Statistics:}
    Probabilistic Models,
    Theoretical Concepts of Machine Learning.

    \item \textbf{Computing Systems:}
    Parallel Computing,
    Optimum and Adaptive Signal Processing Systems.

    \item \textbf{Applications:}
    Computer Vision,
    Introduction to Autonomous Vehicles.

    \item \textbf{Practical Research:}
    Practical Work in AI (Master).

\end{ecvitemize}
}

\ecvitem{Applied Experience}{
\begin{ecvitemize}

    \item Model training, evaluation,
    and hyperparameter optimisation.

    \item Experimental research involving
    deep learning architectures.

    \item Mathematical analysis and
    real-time system modelling.

    \item Statistical analysis and
    data-driven development.

    \item Control-theoretic analysis
    of system stability and reliability.

\end{ecvitemize}
}

% ------------------------------------------------------------
% BACHELOR'S DEGREE
% ------------------------------------------------------------

\ecvtitlelevel{
July 2014 -- June 2019
}{
BSc Mechatronics Engineering
}{
ISCED~6
}

\ecvitem{}{
\href{https://www.ju.edu.jo/}{
University of Jordan
}, Jordan
}

\ecvitem{Core Subjects}{
\begin{ecvitemize}

    \item System Dynamics and Control Theory.

    \item Robotics and Automation.

    \item Embedded Systems and Digital Logic.

    \item Electronics and Programmable Logic Controllers.

    \item Mechanical Engineering and CAD.

    \item Sensors, Actuators, and Measurements.

\end{ecvitemize}
}

\ecvitem{Advanced Topics}{
\begin{ecvitemize}

    \item Microprocessors and Microcontrollers.

    \item Modern and Digital Control Systems.

    \item Signal Processing.

    \item Industrial Process Control.

    \item Power Electronics and Drives.

    \item Electrical Actuators and Motors.

    \item Thermodynamics and Energy Systems.

    \item Measurements and Instrumentation.

\end{ecvitemize}
}

\ecvitem{Final Grade}{
GPA: 3.15
}

\ecvitem{Bachelor's Thesis}{
Design of Intelligent Transportation System
}

\ecvitem{Thesis Description}{
Designed an intelligent transportation system concept
using Vehicle-to-Vehicle (V2V) communication.

The proposed system enables connected vehicles to
cooperatively identify, locate, and track vehicles
without intelligent onboard sensing.

The architecture uses shared information to
construct environmental representations,
estimate potential collisions, and generate warnings.

The research extends cooperative vehicle
communication concepts to improve the
representation of non-connected road users.
}

% ============================================================
% TECHNICAL PROJECTS
% ============================================================

\ecvsection{Selected technical projects}

% ------------------------------------------------------------
% LACOODA PROCESSOR
% ------------------------------------------------------------

\ecvtitle{
2026 -- Present
}{
LACOODA -- Custom Processor Architecture
}

\ecvitem{}{
Independent Hardware Engineering Project
}

\ecvitem{}{
\begin{ecvitemize}

    \item Development of a custom processor
    architecture using SystemVerilog.

    \item Design and verification of a
    64-bit arithmetic logic unit (ALU).

    \item Implementation of arithmetic,
    logical, comparison, and shift operations.

    \item Development of a register file,
    status register, and supporting
    processor components.

    \item Exploration of instruction formats,
    instruction decoding, and processor datapaths.

    \item Development of SystemVerilog
    testbenches for functional verification.

    \item Use of Icarus Verilog,
    Verilator, and GTKWave.

    \item FPGA synthesis and investigation
    of resource utilisation and timing.

\end{ecvitemize}
}

\ecvitem{Technologies}{
SystemVerilog, FPGA, RTL Design,
Digital Logic, Computer Architecture,
Icarus Verilog, Verilator, GTKWave,
Linux, Git
}

% ------------------------------------------------------------
% LACOODA GRAPHICS
% ------------------------------------------------------------

\ecvtitle{
2026 -- Present
}{
LACOODA -- GPU Architecture
}

\ecvitem{}{
Ongoing Hardware Architecture Project
}

\ecvitem{}{
\begin{ecvitemize}

    \item Development of a custom
    graphics accelerator architecture
    using SystemVerilog.

    \item Design exploration for
    a 2D rendering subsystem.

    \item Investigation of command processing,
    graphics registers, memory interfaces,
    and framebuffer organisation.

    \item Integration planning with a custom
    processor and shared memory infrastructure.

    \item Hardware resource and
    FPGA feasibility analysis.

\end{ecvitemize}
}

\ecvitem{Technologies}{
SystemVerilog, GPU Architecture,
Digital Design, FPGA,
Memory Systems, RTL
}

% ------------------------------------------------------------
% AI ACCELERATOR RESEARCH
% ------------------------------------------------------------

\ecvtitle{
2026 -- Present
}{
Machine Learning Accelerator Architecture
}

\ecvitem{}{
Independent Research and Architecture Exploration
}

\ecvitem{}{
\begin{ecvitemize}

    \item Investigation of FPGA-based
    machine learning acceleration.

    \item Exploration of arithmetic and
    computational architectures suitable
    for neural network inference.

    \item Research into efficient implementation
    of mathematical operations in hardware.

    \item Investigation of resource-efficient
    datapaths and computational units.

    \item Exploration of CPU, GPU,
    and TPU integration on resource-constrained FPGAs.

    \item Research interest in using
    large language models to assist
    RTL design and hardware verification.

\end{ecvitemize}
}

\ecvitem{Research Areas}{
AI Accelerators, FPGA,
Hardware-Software Co-design,
Neural Network Inference,
Digital Arithmetic,
LLM-Assisted Hardware Design
}

% ------------------------------------------------------------
% OPENJOEY2
% ------------------------------------------------------------

\ecvtitle{
2026 -- Present
}{
OpenJoey2 -- C++ Game Engine
}

\ecvitem{}{
Independent Software Engineering Project
}

\ecvitem{}{
\begin{ecvitemize}

    \item Development of a custom
    game engine using C++17.

    \item Design of modular systems
    for game-state representation
    and rule execution.

    \item Development of data structures
    for cards, zones, players,
    and game-state transitions.

    \item Exploration of instruction-based
    execution of game rules and effects.

    \item Development of semantic
    representations of card effects.

    \item Use of Raylib for
    graphics and user interaction.

    \item Investigation of deterministic
    state transitions and replayable
    execution models.

\end{ecvitemize}
}

\ecvitem{Technologies}{
C++17, CMake, Raylib,
Algorithms, Data Structures,
State Machines, JSON,
Software Architecture
}

% ============================================================
% PERSONAL SKILLS
% ============================================================

\ecvsection{Personal skills}

% ------------------------------------------------------------
% LANGUAGE SKILLS
% ------------------------------------------------------------

% Language proficiency values are left unchanged from
% the source CV, where they were commented out.
% Enable only after confirming the current proficiency.

% \ecvmothertongue{Arabic}
% \ecvlanguageheader
% \ecvlanguage{English}{B2}{B2}{B2}{B2}{B2}
% \ecvlanguage{German}{A1}{A1}{A1}{A1}{A1}
% \ecvlanguagefooter

% ------------------------------------------------------------
% COMPUTER SKILLS
% ------------------------------------------------------------

\ecvblueitem{Computer skills}{
\begin{ecvitemize}

    \item \textbf{Programming Languages:}
    C++, C, Python, MATLAB,
    SystemVerilog, Verilog.

    \item \textbf{Artificial Intelligence:}
    Machine Learning, Deep Learning,
    PyTorch, Neural Networks,
    LSTM, Reinforcement Learning,
    Probabilistic Models.

    \item \textbf{Computer Vision:}
    OpenCV, Image Processing,
    Computer Vision Algorithms,
    Data Analysis.

    \item \textbf{Digital Hardware:}
    RTL Design, Digital Logic,
    Combinational and Sequential Circuits,
    ALU Architecture, Register Files,
    Processor Datapaths.

    \item \textbf{FPGA Development:}
    FPGA Architecture,
    RTL Simulation,
    Functional Verification,
    Synthesis,
    Resource Utilisation Analysis,
    Timing Analysis.

    \item \textbf{Hardware Verification:}
    SystemVerilog Testbenches,
    Icarus Verilog,
    Verilator,
    GTKWave.

    \item \textbf{Hardware Research:}
    Processor Architecture,
    GPU Architecture,
    AI Accelerators,
    Hardware-Software Co-design.

    \item \textbf{Systems and Control:}
    Control Theory, Simulink,
    Embedded Systems,
    PLC, Sensors,
    Signal Processing.

    \item \textbf{Scientific Computing:}
    Statistical Analysis,
    Mathematical Modelling,
    Numerical Computing,
    Parallel Computing.

    \item \textbf{Software Engineering:}
    Object-Oriented Programming,
    Algorithms, Data Structures,
    CMake, Qt, Raylib.

    \item \textbf{Infrastructure and Tools:}
    Linux, Bash, Git,
    VS Code, Selenium.

    \item \textbf{Additional Technologies:}
    OpenCL, React,
    Large Language Models.

\end{ecvitemize}
}

% ------------------------------------------------------------
% RESEARCH INTERESTS
% ------------------------------------------------------------

\ecvblueitem{Research interests}{
\begin{ecvitemize}

    \item Machine Learning Accelerators.

    \item FPGA-Based Artificial Intelligence.

    \item Efficient Neural Network Inference.

    \item Computer Architecture and Digital Design.

    \item Hardware-Software Co-design.

    \item LLM-Assisted Hardware Development.

    \item Embedded Artificial Intelligence.

    \item Computer Vision and Intelligent Systems.

    \item Signal Processing and Control Systems.

\end{ecvitemize}
}

% ------------------------------------------------------------
% OTHER SKILLS
% ------------------------------------------------------------

\ecvblueitem{Other skills}{
\begin{ecvitemize}

    \item Competitive programming
    and algorithmic problem-solving.

    \item Research-oriented software development.

    \item Mathematical analysis
    and system modelling.

    \item Experience in interdisciplinary
    engineering environments.

    \item Technical documentation
    and scientific communication.

    \item Experience supporting university
    teaching and practical exercises.

\end{ecvitemize}
}

% ============================================================
% ADDITIONAL INFORMATION
% ============================================================

\ecvsection{Additional information}

% ------------------------------------------------------------
% COMPETITIVE PROGRAMMING
% ------------------------------------------------------------

\ecvblueitem{
Competitive Programming Achievements
}{
Participated in national, regional,
and international programming competitions
in algorithms, problem-solving,
and data structures.

C++ was the primary programming language.

\begin{ecvitemize}

    \item \textbf{IEEEXtreme 2019:}
    93rd out of approximately 9,400 participants.

    \item \textbf{Africa and Arab Collegiate
    Programming Contest 2019:}
    28th place.

    \item \textbf{Jordan Collegiate
    Programming Contest 2019:}
    5th place.

    \item \textbf{Jordan Collegiate
    Programming Contest 2018:}
    13th place.

    \item \textbf{Arab Collegiate
    Programming Contest 2017:}
    92nd place.

\end{ecvitemize}

Certificates:

\href{
https://github.com/shamekhing/Certificates
}{
github.com/shamekhing/Certificates
}
}

% ------------------------------------------------------------
% PROFESSIONAL SUMMARY
% ------------------------------------------------------------

\ecvbigitem{Professional summary}{
AI and software engineer combining practical
research experience with a multidisciplinary
foundation in computer science,
mechatronics, and control engineering.

Professional experience includes applied
machine learning, research engineering,
scientific software, industrial software
development, and automation.

Independent technical projects extend
this experience into computer architecture,
SystemVerilog, FPGA development,
and hardware acceleration.

Interested in research and engineering
positions involving artificial intelligence,
embedded computing, digital hardware,
and efficient computational systems.
}

% ============================================================
% END DOCUMENT
% ============================================================

\end{europasscv}

\end{document}